\section{Robustness}\label{sec:robust}
% =============================================================================================
% E5 (experiments/e5_robustness.py; parts grid, split, simpower as run; gridsize, gridpower,
% gridrealvol not run) and the post-hoc architecture ablation (experiments/e5_architecture.py ->
% results/e5arch/ -> tables e5_arch_sim.tex, e5_arch_real.tex).
% Main text: tab:manual-e5-sensitivity, fig:e5-sensitivity, tab:manual-e5-split (condensed tables
% written by paper/scripts/e5_tables.py; the generated e5_sensitivity.tex and e5_split.tex are too
% wide for the main text and are in app:robust as tab:e5-sensitivity, tab:e5-split).
% Numbers: paper/tables/e5_*.tex, manual_e5_*.tex, results/e5/*.csv, results/e5arch/*.csv,
% results/e3/e3_{power,realvol}_{raw,summary}.csv. Derived counts (orderings, McNemar tests):
% recomputed from the raw CSVs. Paired alternative/boundary outcomes of the rule in D1/D2 at both
% splits (minimum boundary pmed, alternatives below it): paper/scripts/e5_split_discrimination.py
% -> paper/tables/manual_e5_split_paired.tex (input in app:robust as tab:manual-e5-split-paired).
% Architecture pilot: docs/DESIGN.md change log 2026-11-16 ("four E2 designs at n = 500"; the pilot
% itself is not in the repository).
% =============================================================================================

E5 reruns the real-data decision rule of \cref{sec:method-decision} under one-at-a-time changes of
its tuning choices and under an alternative \TRAIN/\TEST{} split, and examines the optimiser
settings of the critic on simulated data; a post-hoc ablation replaces the critic network. Every
real-data run uses the 100 tickers, fitted models, five synthetic seeds and test randomness of
E3, so the default configuration reproduces E3 exactly: its 50 p-values (absolute \WC{} and
anchored tests, five generators, five seeds) are identical to those of E3. Three planned parts
were not run for lack of compute (\cref{app:robust}): the size of the absolute test under each
configuration, and the power and boundary size of the rule under each configuration in designs
D1 and D2. A non-default configuration therefore has no power estimate; its verdict can be \BAD{}
or \INCONCL{} but not \GOOD, and \INCONCL{} there only means that the anchored test did not
reject. The size of the default pipeline is measured in E1 (\cref{tab:e1-size-sim}): at nominal
5\% the absolute \WC{} and anchored tests reject a true null at 0.085 and 0.080 when the synthetic
panel is a joint draw, but at 0.270 and 0.155 when the synthetic tickers are independent, as they
are for the univariate generators examined here.

\paragraph{Sensitivity grid.}
The grid changes one setting at a time (\cref{tab:manual-e5-sensitivity,fig:e5-sensitivity}): the
window length $L\in\{5,10,60\}$ (default 20), the width and depth of the critic ($32\times2$ and
$256\times4$; default $128\times3$), the number of training steps (500 and 3000; default 1000),
the learning rate ($10^{-4}$; default $10^{-3}$), the gradient-penalty weight ($\lgp=1$; default
10), the number of time blocks ($K=4$; default 10), the training fraction (0.3 and 0.7; default
0.5) and the HAC lag rule (the rule of \citet{newey1994automatic} instead of the plug-in of
\citet{andrews1991heteroskedasticity}). Each change applies to the absolute and to the anchored
test. Three results hold in every configuration. First, the absolute \WC{} test rejects at 5\% in
69 of the 70 configuration--generator cells; the exception is the stationary bootstrap with $K=4$
($\pmed=0.095$). Second, the anchored test never rejects \GARCHt{}: its $\pmed$ ranges from 0.100
($L=5$) to 1.000. Third, the extremes of the ranking of the likelihood models do not change. In
every configuration GBM has the smallest anchored $\pmed$ and \GARCHt{} the largest (shared with
Heston at 1.000 when the training fraction is 0.3), and the median absolute statistic decreases
from GBM to Merton, Heston and \GARCHt{}. The anchored order of Merton and Heston is not stable:
Heston has the smaller $\pmed$ with $L=60$ (0.070 against 0.110) and with the $256\times4$ network
(0.028 against 0.029). Both statistics thus rank GBM below Merton and Merton below \GARCHt, the
order of the predictive log-score (\cref{tab:e3-agreement}).

The \BAD{} verdicts are less stable. The anchored test rejects GBM in 13 of the 14
configurations, Merton in 11 and Heston in 9. All nine changes against the default go from
rejection to non-rejection; no configuration rejects a generator that the default does not
reject. They occur where the anchored test has fewer evaluation dates or a weaker critic: longer
windows ($L=60$ leaves 10 windows per ticker in \TEST; Merton $\pmed=0.110$, Heston 0.070), a
30\% training share (GBM 0.170, Merton 0.671, Heston 1.000), four time blocks (Merton 0.236,
Heston 0.379) and, for Heston only, the small network (0.102) and the lower learning rate
(0.079). Shorter windows, the larger network, fewer or more steps, $\lgp=1$, a 70\% training
share and the Newey--West rule change no decision, but the narrowest decision of the grid is
among them: with 500 steps Heston has $\pmed=0.0499$, printed as 0.050 in
\cref{tab:manual-e5-sensitivity} and just below the level. At the default, Merton ($\pmed=0.005$) and
Heston (0.017) are closer to the level than GBM ($<0.001$). \Cref{fig:e5-sensitivity} also shows
why the seeds are combined: the single-seed p-values of one generator in one configuration
spread over several orders of magnitude (\cref{rem:combine-seeds}).

\begin{table}[tbp]
  \centering
  \caption{Sensitivity of the E3 decisions to one-at-a-time changes of the pipeline (E5 grid;
    100 tickers, \TEST{} 2023-01-03 to 2025-06-16, five synthetic seeds, $\alpha=0.05$).}
  \label{tab:manual-e5-sensitivity}
  \small
  \fittable{manual_e5_sensitivity}
  \tablenote{Generator columns: $\pmed=\min\{1,2\med_sp_s\}$ over the five seeds of the
    Driscoll--Kraay p-values (time-block split) of the anchored relative test, followed by the
    verdict (B = \BAD, I = \INCONCL). $^\dagger$: the decision at 5\% differs from the default row.
    $^\ddagger$: relative test invalid for the stationary bootstrap (no verdict). Last column:
    number of the five generators that the absolute \WC{} test ($\pmed$ over the same seeds)
    rejects at 5\%. The power check was not rerun for any configuration, so a non-default verdict
    is \BAD{} or \INCONCL{}; the default verdicts coincide with \cref{tab:e3-verdicts}, whose power
    comes from \cref{tab:e3-power}. All p-values: \cref{tab:e5-sensitivity} in \cref{app:robust}.}
\end{table}

\begin{figure}[tbp]
  \centering
  \inputfigure{e5_sensitivity}
  \caption{p-values of the E5 grid (100 tickers, \TEST{} 2023-01-03 to 2025-06-16, five synthetic
    seeds). Left: anchored relative test; right: absolute \WC{} test; one row per configuration.
    Open markers: the five single-seed Driscoll--Kraay p-values; filled markers:
    $\pmed=\min\{1,2\med_sp_s\}$. Dashed line: $\alpha=0.05$. p-values below $10^{-4}$ are drawn at
    the left edge. The stationary bootstrap is faint in the left panel because its relative test
    is invalid.}
  \label{fig:e5-sensitivity}
\end{figure}

\paragraph{Alternative split.}
With \TRAIN{} returns from 2018-01-03 to 2021-12-31 and \TEST{} 2022-01-03 to 2025-06-16 (50 and 43 windows
per ticker), the five generators are refitted and the full E3 battery is rerun with five seeds;
the power and boundary size of the rule are estimated at the sizes of this split with 25
replications per design (\cref{tab:manual-e5-split}). The absolute tests reject fewer
generators, although \TEST{} is longer than at the main split. The \WC{} test rejects only GBM
($\pmed<0.001$; Merton 0.055, Heston 0.076, \GARCHt{} 1.000), against all five generators at the
main split, and no valid window test rejects \GARCHt{} or the bootstrap. The anchored test rejects
no generator (GBM 0.091, Merton 0.704, \GARCHt{} 0.516, Heston 0.966). The power of the rule is
1.00 [0.87, 1.00] in D1 but 0.60 [0.41, 0.77] in D2, below $\bmin$, so every verdict is
\INCONCL{}. None of the three \BAD{} verdicts of the main split is reproduced.

The boundary sizes show that \INCONCL{} is the right outcome, for a reason that the rule does not
encode. In D1 the rule never rejects at the boundary (0.00 [0.00, 0.14], Clopper--Pearson). In D2
its boundary size is 0.60 [0.39, 0.79]: a generator that reproduces the pooled \TRAIN{} window law
exactly is declared worse than history in 15 of 25 replications. These are the same 15
replications in which the alternative $G^*=1.2\,G^0$ is rejected, and the alternative is rejected
in no other replication (\cref{app:robust}). At the nominal threshold the D2 power of the rule is
therefore not interpretable. The critic still separates the two laws: in each of the 15 rejecting
replications the alternative has $\pmed\le7.8\times10^{-6}$, below the smallest boundary $\pmed$ of
the 25 replications ($3.8\times10^{-5}$), so a critical value just below that smallest boundary
value, itself estimated from these replications, would reject the boundary in none of the 25 and
the alternative in the same 15 (\cref{tab:manual-e5-split-paired} in \cref{app:robust}). What fails at this split is the calibration of the nominal
threshold, not discrimination; whether the nominal test fires depends on the draw of the
replication. The median anchored statistic at the D2 boundary is 0.231, against 0.012 at the main
split, so the anchored target at the pooled-\TRAIN{} boundary is far from zero, the case that
\cref{rem:nonstat-past} anticipates for a heterogeneous past. Three consequences follow. First, a
\BAD{} verdict at this split, read at the nominal threshold, would have said nothing about the
generator. Second, the rule \eqref{eq:verdict} checks power but not boundary size, and
its \BAD{} branch can be interpreted only where the estimated boundary size is close to $\alpha$,
as at the main split (D1 0.00 [0.00, 0.04], D2 0.08 [0.04, 0.15]; \cref{tab:e3-power}). A rule
that turns \BAD{} into \INCONCL{} whenever the Clopper--Pearson interval of the boundary size lies
above $\alpha$ would encode this; it is not the rule that produced our tables. Third, under
non-stationarity a verdict is a statement about one pair of periods: the past that serves as the
benchmark changes with the split, and so do the verdicts and the absolute rejections.

\begin{table}[tbp]
  \centering
  \caption{Alternative split (E5): \TRAIN{} 2018-01-03 to 2021-12-31, \TEST{} 2022-01-03 to
    2025-06-16, same 100 tickers, generators refitted; five synthetic seeds; rule power and
    boundary size from 25 replications per design; $\alpha=0.05$.}
  \label{tab:manual-e5-split}
  \small
  \fittable{manual_e5_split}
  \tablenote{Entries are $\pmed$ over the five seeds of the Driscoll--Kraay p-values (time-block
    split) of the absolute \WC{} test and of the anchored relative test; bold: absolute rejection
    at 5\%. $^\ddagger$: relative test invalid for the stationary bootstrap (verdict n/a). Rule
    power (synthetic volatility $\times1.2$) and boundary size: share of replications (number in
    parentheses) in which $\pmed\le0.05$ over five seeds, at the sizes of each split, with Wilson
    (power) and Clopper--Pearson (boundary size) 95\% intervals; main split as in
    \cref{tab:e3-power}. All tests of the battery: \cref{tab:e5-split} in \cref{app:robust}.}
\end{table}

\paragraph{Optimiser settings.}
On simulated independent windows with exact permutation p-values ($n=500$ windows per sample, 50
replications per cell; \cref{tab:e5-simpower}), the default critic has power 0.40 [0.28, 0.54]
against a 10\% increase in the volatility of Student-$t_5$ returns and 0.30 [0.19, 0.44] against
GARCH persistence 0.90 instead of 0.98. With learning rate $10^{-4}$ and the default 1000 steps,
the power against the volatility shift rises to 0.78 [0.65, 0.87], while against clustering it
is 0.34 [0.22, 0.48], not detectably different. Power is not monotone in the training budget: at
rate $10^{-4}$ it is highest at 1000 steps, and 500 steps under-train the critic. The level does
not depend on any of these settings (\cref{prop:exact-level}); the power does. The defaults were
fixed before any E2 run and were not re-tuned on these results, so the \WC{} power of
\cref{sec:sim} is that of the default and may understate what a differently trained critic
achieves against volatility-level deviations. On real data the lower rate did not reject more;
it turned Heston from \BAD{} to \INCONCL{} ($\pmed=0.079$). Its rule power in D1 and D2 was not
estimated, so this change can be attributed neither to size nor to power, and whether the lower
rate would raise the D2 power of the rule is unknown.

\paragraph{Critic architecture.}
The network matters at least as much as the optimiser, but the evidence for the alternative
network is weaker than for the default. A time-pooled convolutional critic, which averages one
learned function of nine-day patches over the window (\cref{app:robust}), was chosen after the
real-data results were known and after a pilot on four E2 designs at $n=500$ had shown it to be
more powerful than the perceptron; the pilot is not part of the reported experiments. It is
therefore a post-hoc ablation, and the pre-specified verdicts remain those of the perceptron. In
replications paired with those of E2 and E2-dep (\cref{tab:e5-arch}), it keeps the level with
independent windows (0.06 and 0.05 in 200 null replications) and raises the power in every
alternative design, significantly in all but the block design of E2-dep: from 0.10 to 0.98
against lighter tails ($t_{10}$ against $t_4$), from 0.10 to 0.92 when every synthetic window is
GBM instead of Merton, and from 0.36 to 0.88 against persistence 0.90. Because the network was
selected on E2 designs of this kind, possibly on some of these very cells, the E2 gains are
optimistic. The two E2-dep cells ($n=2000$), outside the pilot as recorded, are the closest to
held-out evidence: 0.90 and 0.96 against 0.46 and 0.86 for the perceptron (exact McNemar
$p<0.001$ and $p=0.18$). In the designs on which the pooled critic was selected, the weakness of
the \WC{} test against marginal deviations in \cref{sec:sim} is thus largely a property of the
multilayer perceptron. The size of the pooled critic under the dependent protocol (E1(a)) was not
measured. On the S\&P~500 panel it gives the same five decisions (\cref{tab:e5-arch-real}). Its
rule power and boundary size were estimated only in the first 20 replications of D1 and D2, where
they equal those of the perceptron on the same replications: power 1.00 and 0.45 [0.26, 0.66],
boundary size 0.00 and, in D2, 0.20 [0.06, 0.44] (Clopper--Pearson), against 0.08 [0.04, 0.15]
over all 100 replications of the perceptron. These replications are not representative of E3 at
the D2 boundary and cannot separate the two networks. Whether a more powerful critic, or the
lower learning rate, would lift the D2 power above $\bmin$ is therefore open.

\paragraph{What is robust.}
At the main split, across the thirteen changes of the grid, the absolute \WC{} test rejects every
parametric generator (and the bootstrap in every configuration but the one with four time blocks,
$\pmed=0.095$), the anchored test does not reject \GARCHt{}, and both statistics rank GBM below
Merton and Merton below \GARCHt. At the alternative split the absolute \WC{} test rejects only GBM
and the anchored test rejects no generator. \GARCHt{} is therefore never \BAD{}; it is never
\GOOD{} either, because its D2 power is below $\bmin$ at both splits and was not estimated in the
grid. GBM is \BAD{} in all but one configuration of the grid. Four conclusions are not robust. The
\BAD{} verdicts of Merton and Heston disappear under three and five of the thirteen changes, and
that of Heston with 500 steps rests on $\pmed=0.0499$. No \BAD{} verdict survives the change of
split, and at the alternative split the nominal threshold of the rule is not calibrated at the
boundary in the design with the real volatility path, although the critic still separates the
alternative from the boundary. The absolute rejections of Merton, Heston, \GARCHt{} and the
bootstrap also disappear at that split. The power of the \WC{} test depends on the optimiser
settings, not monotonically, and on the network. The rule power on real data, on which a \GOOD{}
verdict depends, has been estimated only for the default optimiser, with the perceptron and, in
20 replications, with the pooled critic. The exact level for independent windows holds for every
setting. On real data, calibration at the relative boundary is a
property of the chosen periods and has to be checked for each split, as the boundary size does.
Whether a generator is good enough for a task such as portfolio construction has to be checked
on that task (\cref{sec:markowitz}).
